\documentclass[11pt]{report}

\usepackage[margin=1in]{geometry}
\usepackage[T1]{fontenc}
\usepackage[utf8]{inputenc}
\usepackage{lmodern}
\usepackage{microtype}
\usepackage{amsmath,amssymb,amsthm,mathtools}
\usepackage{array}
\usepackage{booktabs}
\usepackage{longtable}
\usepackage{enumitem}
\usepackage{xcolor}
\usepackage{fancyhdr}
\usepackage{hyperref}
\usepackage{listings}
\usepackage{stmaryrd}
\usepackage{tikz}
\usepackage{tikz-cd}
\usepackage{algorithm}
\usepackage{algpseudocode}

\hypersetup{colorlinks=true,linkcolor=blue!60!black,urlcolor=blue!60!black,citecolor=blue!60!black}
\usetikzlibrary{arrows.meta,positioning,calc,fit}
\definecolor{chapterblue}{HTML}{1F4E79}
\definecolor{accentgold}{HTML}{C69214}

\lstdefinestyle{deppy}{
  language=Python,
  basicstyle=\ttfamily\small,
  keywordstyle=\color{blue!60!black},
  commentstyle=\color{green!40!black},
  stringstyle=\color{red!50!black},
  showstringspaces=false,
  frame=single,
  breaklines=true,
  columns=fullflexible,
  mathescape=true
}
\lstset{style=deppy}

\newcommand{\DepPy}{\textsc{deppy}}
\newcommand{\Site}{\ensuremath{\mathsf{Site}}}
\newcommand{\Layer}{\ensuremath{\mathsf{Layer}}}
\newcommand{\Sem}{\ensuremath{\mathsf{Sem}}}
\newcommand{\Type}{\ensuremath{\mathsf{Type}}}
\newcommand{\Trust}{\ensuremath{\mathsf{Trust}}}
\newcommand{\Core}{\ensuremath{\mathsf{Core}}}
\newcommand{\HybridTy}{\ensuremath{\mathsf{HybridTy}}}
\newcommand{\py}[1]{\texttt{\detokenize{#1}}}
\newcommand{\Hc}[1]{\check{H}^{#1}}
\newcommand{\Cech}{\check{C}}
\newcommand{\op}{\mathrm{op}}
\newcommand{\TO}{\ensuremath{T_O}}
\newcommand{\Bool}{\ensuremath{\mathbb{2}}}
\newcommand{\Score}{\ensuremath{\sigma}}
\newcommand{\rank}{\operatorname{rank}}
\newcommand{\im}{\operatorname{im}}
\newcommand{\Hom}{\operatorname{Hom}}
\newcommand{\Id}{\operatorname{Id}}
\DeclareRobustCommand{\repo}[1]{\path{#1}}
\emergencystretch=2em
\linespread{1.03}
\setlength{\headheight}{14pt}
\pagestyle{fancy}
\fancyhf{}
\fancyhead[LE,RO]{\thepage}
\fancyhead[LO]{\nouppercase{\rightmark}}
\fancyhead[RE]{\nouppercase{\leftmark}}
\fancyfoot[C]{\textcolor{chapterblue}{\DepPy{} --- cohomological typing}}
\renewcommand{\headrulewidth}{0.4pt}
\renewcommand{\footrulewidth}{0pt}

\theoremstyle{definition}
\newtheorem{definition}{Definition}[chapter]
\newtheorem{example}[definition]{Example}
\newtheorem{remark}[definition]{Remark}
\newtheorem{assumption}[definition]{Assumption}

\theoremstyle{plain}
\newtheorem{theorem}[definition]{Theorem}
\newtheorem{lemma}[definition]{Lemma}
\newtheorem{proposition}[definition]{Proposition}
\newtheorem{corollary}[definition]{Corollary}

\title{Cohomological Typing for Absolutely Everything\\
\large A \DepPy{} Monograph on Sheaf-Descent, Hybrid Dependent Types, Trust, and Generation}
\author{Repository-integrated draft}
\date{\today}

\begin{document}
\begin{titlepage}
\thispagestyle{empty}
\centering
\vspace*{2.5cm}
{\color{chapterblue}\rule{0.82\textwidth}{1.2pt}\par}
\vspace{1.25cm}
{\Huge\bfseries Cohomological Typing for Absolutely Everything\par}
\vspace{0.6cm}
{\Large\color{chapterblue}A \DepPy{} monograph on sheaf-descent, hybrid dependent types, trust, and generation\par}
\vspace{1.3cm}
{\large Repository-integrated draft\par}
\vspace{0.25cm}
{\large \today\par}
\vfill
\begin{minipage}{0.84\textwidth}
\centering
\textit{Programs are covered by local observations; types are the sections attached to those observations; proofs, semantic judgments, and generated artifacts succeed when the sections glue.}
\end{minipage}
\vfill
{\color{accentgold}\rule{0.82\textwidth}{1.2pt}\par}
\end{titlepage}

\pagenumbering{roman}

\begin{abstract}
This monograph takes the strongest idea running through \DepPy{} and pushes it to its
logical limit: if typing is fundamentally a local-to-global problem, then the natural
language for typing is not a flat set of judgments but a sheaf-like organization of local
sections, overlap maps, trust labels, proof objects, and executable artifacts. The
repository already contains several versions of this idea. The foundational
\repo{src/deppy/types/} hierarchy presents a sheaf-descent semantic typing substrate for
Python. The \repo{src/deppy/hybrid/} tree extends that substrate into a layered object
whose coordinates are INTENT, STRUCTURAL, SEMANTIC, PROOF, and CODE, equipped with a
trust lattice and an oracle monad. The equivalence machinery, anti-hallucination
pipeline, theory packs, type expansion modules, zero-change analyzers, and prompt-driven
agent all reuse the same local-to-global intuition.

The goal of this document is to describe that entire ecosystem as one research program:
\emph{cohomological typing for absolutely everything}. Here ``absolutely everything''
means not that every software question reduces to a single trivial theorem, but that the
same geometric and dependent-typing viewpoint can organize foundational types,
verification obligations, solver fragments, semantic oracle judgments, library theories,
interprocedural transport, program equivalence, hallucination detection, prompt-driven
planning, and even domain ideation. The manuscript therefore reads the repository as a
single layered system rather than as a bag of isolated prototypes.

The central claims are straightforward. First, the plain type system in
\repo{src/deppy/types/} is still the object language: it gives \DepPy{} its carrier
hierarchy, refinement types, dependent functions and pairs, equality-like objects,
indexed families, and subtype obligations. Second, the hybrid system does not replace
that language; it lifts it into a richer semantic object
\[
  \HybridTy : (\Site(P) \times \Layer)^{\op} \to \Type \times \Trust,
\]
where coherence is measured by the vanishing of first \v{C}ech cohomology and semantic
predicates live inside the oracle monad
\[
  \TO(A) = A \times \mathsf{TrustLevel} \times \mathsf{Confidence} \times \mathsf{Provenance}.
\]
Third, many subsystems that appear separate in practice become legible as specializations
of the same architecture: bugs as non-gluing cocycles, equivalence as gluable local
isomorphism, theory packs as local solver sheaves, mixed-mode syntax as a user-facing
bridge into plain obligations and hybrid contracts, and the generation agent as a typed
cover-growing machine.

The document is deliberately repository-specific. It names the relevant files,
acknowledges the places where the repository duplicates concepts or shifts vocabulary,
and treats current examples---including trading, zero-change verification, and
anti-hallucination---as evidence for a broader thesis rather than as disconnected demos.
\end{abstract}

\tableofcontents
\clearpage
\pagenumbering{arabic}

\part{Foundations}
\chapter{Introduction}

\DepPy{} is easiest to misunderstand when its different subsystems are read as separate
one-off experiments: a plain type hierarchy here, a hybrid dependent type story there, a
cohomological bug-finding paper somewhere else, and an agent-driven generation runtime on
top. The working thesis of this manuscript is that those pieces are better read as one
architecture. The architecture begins from a simple observation: software understanding is
local-to-global. We learn facts at boundaries, guards, call sites, proof obligations,
library interfaces, and traces. We then try to compose those local facts into a global
judgment about safety, meaning, equivalence, or implementation quality. When local facts
agree, we obtain a global section. When they fail to agree, we do not merely have a vague
error---we have a gluing obstruction.

That is why cohomological language keeps reappearing across the repository. In the older
sheaf-oriented work, a program is analyzed through a site category of observation points,
a semantic presheaf of local facts, and the computation of \v{C}ech cohomology to locate
where those facts fail to glue. In the hybrid work, the same local-to-global idea is
lifted to a five-layer product site where intent, structural facts, semantic judgments,
proof artifacts, and code must agree. In the equivalence machinery, one asks whether
local isomorphisms glue to a global one. In anti-hallucination, one asks whether a
candidate artifact inhabits the expected type on both its structural and semantic faces.
In prompt-driven generation, one grows a cover of typed module obligations and tries to
glue them into a coherent project.

The phrase ``cohomological typing for absolutely everything'' is intentionally ambitious,
but it is not meant as parody. It means that the same mathematical stance can organize
far more of the repository than ordinary type-checking alone. Typing becomes a universal
language of structured local evidence: value-level facts, control-flow facts, library
facts, semantic judgments, proof obligations, and design intent all enter the same story.
The right reaction is not to believe that every workflow literally computes a heavy
cohomology group in the same way. The right reaction is to notice that the repository has
already started treating the local-to-global pattern as its common semantic denominator.
This monograph makes that denominator explicit.

It is worth saying more about what counts as a single project here. A manuscript can become
long without becoming unified. It can merely accumulate examples, slogans, and modules. A
foundational project earns unity only when its later constructions seem to have been latent
in the earlier ones all along. The burden of this monograph is therefore not simply to
explain many parts of the repository. It is to persuade the reader that once one grants the
local-to-global viewpoint, the rest of the repository becomes a chain of increasingly rich
consequences: first exact obligations, then layered evidence, then cohomological diagnosis,
then trust-sensitive repair, then exploratory elaboration, and finally agent-driven
assembly. The paper should feel like a single river that keeps widening, not like a museum
of related specimens.

The two most important distinctions in what follows are these. First, the plain type
system in \repo{src/deppy/types/} is not superseded by the hybrid system; it is the
object language on which the hybrid system depends. It provides the type constructors,
refinement syntax, binder operations, and subtype structure that later layers continue to
use. Second, the hybrid system in \repo{src/deppy/hybrid/} is not merely ``the theorem
prover side'' or ``the LLM side.'' It is a broader extension that adds verification
layers, trust, provenance, semantic predicates, Lean translation, anti-hallucination,
zero-change adoption, and prompt-driven planning. The hybrid story is therefore both more
mathematical and more practical than a simple proof layer.

\paragraph{What this document is trying to do.}
The document has four goals. First, it gives a coherent account of the foundational plain
type system and the layered hybrid extension. Second, it explains how site categories,
presheaves, trust lattices, and dependent formers organize the repository's major ideas.
Third, it shows how apparently distant components---theory packs, renderers, equivalence,
ideation, generation agents, and trading examples---fit into the same architecture.
Fourth, it tries to be honest about tensions in the repository: duplicated trust vocabularies,
heuristic bridges from intent to formalization, and the difference between semantic
acceptance and formal theorem proving.

\paragraph{Why the comparison to a foundational project matters.}
The aspiration behind the present document is not to mimic the vocabulary of better-known
foundational programs, but to borrow one of their virtues: they make the reader feel that a
few basic ideas change the subject all the way down. In that spirit, the early chapters here
will insist on sites, sections, and exact obligations; the middle chapters will insist on
layered evidence, trust, and cohomological obstruction; the later chapters will insist that
generation, theory packs, rendering, and zero-change checking are not external applications
but internal expressions of the same semantics. If the manuscript succeeds, the reader
should feel that the repository is building one universe of discourse for software work, not
merely a bag of clever techniques.

\paragraph{What it means to integrate the repository well.}
To integrate the repository well is not to mention every file exactly once. It is to show
how the files cooperate semantically. A good paper should therefore explain why
\repo{src/deppy/types/refinement.py} matters to \repo{src/deppy/hybrid/mixed_mode/compiler.py},
why \repo{src/deppy/hybrid/core/types.py} matters to \repo{agent/orchestrator.py}, why
\repo{src/deppy/equivalence/} and \repo{src/deppy/render/} are easier to read after one
understands gluing obstructions, and why the zero-change APIs in \repo{deppy.hybrid} are
not merely convenience wrappers but part of a deliberate adoption strategy.

\paragraph{A minimal slogan.}
If a single slogan is needed, it is this:
\begin{quote}
\emph{The plain type system gives \DepPy{} a language of exact semantic obligations.
The hybrid system gives \DepPy{} a language of layered evidence about those obligations.
Cohomology is the bookkeeping device that says whether the layers and local regions really
glue.}
\end{quote}

The remaining chapters expand that slogan until it can carry almost everything the
repository is trying to do.

The reader is therefore invited to hold one question throughout the rest of the manuscript:
\\emph{if this really is one project, what stays invariant as the objects get richer?} The
answer offered by the following chapters is that locality, transport, gluing, obstruction,
repair, and explicit evidence remain invariant. Everything else is the elaboration demanded
by software reality.

\chapter{Scope and Thesis}

This chapter states the central thesis in a way that can absorb both the mathematical and
engineering sides of the repository.

\begin{definition}[Cohomological typing]
Cohomological typing is the discipline of representing typing information as local
sections over a cover of program observations, transporting that information along
restriction maps, and classifying global success or failure by whether those local
sections glue. In the simplest case the data valued in each site is a type-level fact. In
the richer cases used by \DepPy{}, the data also carries trust labels, provenance,
semantic rubrics, proof obligations, and code artifacts.
\end{definition}

The qualifier ``for absolutely everything'' can now be explained more carefully. It does
not mean there is a single theorem that automatically decides all program-analysis
problems. It means the repository repeatedly discovers the same shape of reasoning:
local observations, overlap compatibility, gluing, obstruction, repair, and trust. That
shape is visible in ordinary type checking, in bug diagnosis, in prompt-driven generation,
in semantic oracles, in equivalence testing, in domain-specific theory packs, and in
zero-change adoption for existing Python code. Once the common shape is visible, the
phrase ``typing'' broadens. It comes to mean ``typed local evidence attached to a site and
moved through a gluing discipline.''

The repository pushes this broadening along two axes. The first axis is expressivity. The
plain type system already includes refinements, dependent function and pair formers,
identity-like objects, indexed families, and subtyping. The second axis is evidence.
Hybrid typing extends the foundational language so the same semantic object can be seen at
the level of intent, structural constraints, semantic judgments, proofs, and code. The
result is a repository in which types are no longer merely formulas over values; they are
multi-layer objects that say what is claimed, how it was checked, how much it should be
trusted, and where its remaining ambiguity lives.

A paper with this scope must also be careful about what it is \emph{not} claiming. It is
not claiming that every runtime execution in the repository passes through one uniform
heavy symbolic cohomology engine. It is not claiming that oracle-mediated semantic
judgments have the same epistemic standing as Lean proofs. It is not claiming that every
module is equally mature. Instead, the claim is architectural: the repository is most
coherent when those workflows are viewed as different instances of one local-to-global,
site-aware, evidence-aware semantics.

\paragraph{The central thesis.}
The central thesis of this monograph is that the repository implements, in partial but
already meaningful form, the following picture:
\[
  \Sem : \Site(P)^{\op} \to \mathsf{RefinementLattice}
\]
for the foundational plain system, and
\[
  \HybridTy : (\Site(P) \times \Layer)^{\op} \to \mathsf{RefinementLattice} \times \mathsf{TrustLattice}
\]
for the hybrid system. The former provides the object language of obligations; the latter
provides the layered evidence model in which those obligations are planned, checked,
rendered, translated, discharged, and assembled.

\paragraph{Three design commitments.}
The thesis comes with three commitments that are visible throughout the code base.
\begin{enumerate}[leftmargin=*, itemsep=0.5em]
  \item \textbf{Native Python first.} The mixed-mode surface is designed so users can stay
  inside ordinary Python syntax and still benefit from rich semantics. Surface forms
  degrade gracefully before compilation.
  \item \textbf{Structured evidence, not silent collapse.} Unknown or heuristic facts should
  remain visible as residual or lower-trust data. Trust levels and provenance are meant to
  prevent false certainty.
  \item \textbf{Local-to-global semantics everywhere.} Bugs, equivalence, hallucination,
  library reasoning, and module planning are all interpreted as local facts that may or
  may not glue.
\end{enumerate}

In short: the scope is large, but the architectural thesis is clean.

\chapter{The Cohomological Type Stack}

A mature foundational project feels unified when its central ideas are not introduced as
separate tricks for separate subproblems, but as a small vocabulary that keeps reappearing
until the reader realizes that the whole enterprise is one object seen under several
illuminations. That is one reason some mathematical projects feel inevitable rather than
accidental. They persuade the reader that once a few basic constructions are granted, the
later chapters are not bolt-ons but consequences. \DepPy{} should be read in that spirit.

The repository is easiest to trivialize when its components are described in isolation:
there is a plain type system; there is a hybrid extension; there are trust levels; there is
an oracle monad; there is a cohomological bug story; there are theory packs; there is a
zero-change verifier; there is a generation agent. Read that way, the repository looks like
a mixed research notebook. Read more carefully, it looks like a single stack whose levels
share one grammar. The grammar is the grammar of local sections, transport, gluing,
obstruction, repair, and eventual assembly.

\section*{The one-stack principle}
The most useful way to say this is the \emph{one-stack principle}:
\begin{quote}
\emph{Every major subsystem in \DepPy{} is a manifestation of one cohomological type stack.
The lower levels supply the objects and local obligations; the middle levels supply the
witnesses, trust annotations, and proof boundaries; the upper levels supply the methods by
which typed local material is searched, transported, rendered, and assembled.}
\end{quote}

The stack is not a metaphor pasted over arbitrary code. It is a practical account of how
the repository is organized.

\begin{definition}[The cohomological type stack]
The cohomological type stack is the five-level organization
\[
  \mathsf{Foundation} \to \mathsf{Object} \to \mathsf{Evidence} \to \mathsf{Semantics} \to \mathsf{Assembly}
\]
in which:
\begin{enumerate}[leftmargin=2em, itemsep=0.35em]
  \item \textbf{Foundation} supplies sites, covers, restriction maps, and presheaf structure.
  \item \textbf{Object} supplies the exact type language of carriers, refinements, and
  dependent formers.
  \item \textbf{Evidence} lifts those objects into hybrid witnesses distributed across
  intent, structure, semantics, proof, and code.
  \item \textbf{Semantics} classifies and composes evidence via trust lattices, oracle
  effects, theorem translation, and residual obligations.
  \item \textbf{Assembly} grows, repairs, and glues local artifacts into runnable or
  checkable systems.
\end{enumerate}
\end{definition}

The reward for naming the stack explicitly is that vocabulary starts to stabilize. A site is
a place where one may speak locally. A section is typed local evidence. Transport is the
movement of that evidence along a semantic map. Gluing is the passage from local evidence to
a global artifact. Obstruction is non-gluing. Repair is a refinement of the cover, the local
claim, or the evidence attached to the claim. The same words remain meaningful whether one is
speaking about subtyping, hallucination, theory-pack discharge, or multi-module synthesis.

\section*{Five manifestations of one vocabulary}
The repository's apparent diversity can now be restated as a controlled repetition of this
vocabulary.

At the foundational level, one has sites and presheaves. A function is not seen as one flat
blob of semantics but as a cover of boundaries, guards, intermediate values, call sites, and
error-sensitive observations. That is already a local-to-global stance.

At the object level, one has the plain language in \repo{src/deppy/types/}. Here the
repository defines the actual symbolic objects that higher levels will manipulate:
\py{TypeBase}, refinement types, $\Pi$-types, $\Sigma$-types, equality-like constructions,
indexed families, and subtype obligations. This level says what must be true.

At the evidence level, one has the hybrid language in \repo{src/deppy/hybrid/core/}. Here
those same obligations acquire multiple manifestations: intent text, structural predicates,
semantic rubrics, proof artifacts, and code. This level says how a claim is witnessed and
how much confidence should be placed in the witness.

At the semantic level, one has trust, oracle effects, proof translation, discharge, and the
bookkeeping of residual honesty. This level says what kind of evidence one has, what it is
worth, and how it composes.

At the assembly level, one has type expansion, theory packs, zero-change checking, and the
generation agent. This level says how local pieces are found, prioritized, and glued into a
coherent system.

The central argument of this chapter is that these are not five separate research projects.
They are five manifestations of one project.

\begin{figure}[h]
\centering
\begin{tikzpicture}[
  every node/.style={align=center},
  box/.style={draw=chapterblue, rounded corners=2pt, fill=chapterblue!6, minimum width=10.6cm, minimum height=0.95cm},
  note/.style={draw=accentgold, rounded corners=2pt, fill=accentgold!12, minimum width=3.0cm, minimum height=0.8cm},
  >=Latex
]
\node[box] (foundation) {Foundation \quad sites, covers, restrictions, presheaves};
\node[box, below=0.5cm of foundation] (object) {Object \quad \repo{src/deppy/types/}: carriers, refinements, dependent formers, subtype obligations};
\node[box, below=0.5cm of object] (evidence) {Evidence \quad \repo{src/deppy/hybrid/core/}: layered witnesses across intent, structure, semantics, proof, and code};
\node[box, below=0.5cm of evidence] (semantics) {Semantics \quad trust lattices, oracle monad, Lean translation, discharge, residuals};
\node[box, below=0.5cm of semantics] (assembly) {Assembly \quad theory packs, zero-change analysis, ideation, type expansion, agent synthesis};
\node[note, right=1.2cm of evidence] (glue) {same vocabulary \ section / transport / glue / obstruct / repair};
\foreach \a/\b in {foundation/object,object/evidence,evidence/semantics,semantics/assembly} {
  \draw[->, thick, chapterblue] (\a) -- (\b);
}
\draw[->, thick, accentgold] (object.east) -- (glue.west);
\draw[->, thick, accentgold] (assembly.east) .. controls +(1.2,0.1) and +(0.0,-0.9) .. (glue.south);
\end{tikzpicture}
\caption{The cohomological type stack. The project becomes conceptually unified once the same local-to-global vocabulary is seen to govern each level.}
\end{figure}

\section*{Why the plain and hybrid systems belong to the same project}
The plain and hybrid systems are sometimes described as though they belonged to different
ages of the repository. This is historically understandable and mathematically misleading.
The plain system is not an abandoned first draft; it is the exact object language that makes
the hybrid system possible. The hybrid system is not a repudiation of the plain system; it
is the evidence-bearing lift of the same obligations into a larger universe.

That relation can be stated succinctly. The plain system determines what a value, function,
refinement, or witness \emph{is}. The hybrid system determines what it means to possess
multiple, potentially nonidentical, forms of evidence about that object and to compare those
forms without erasing their differences. The hybrid system therefore presupposes the plain
one in much the way a semantics of paths presupposes points.

\section*{One lexicon across engineering concerns}
Once the stack is visible, many engineering decisions in the repository become easier to
justify.

The mixed-mode syntax is no longer an ergonomic front-end awkwardly bolted onto a verifier.
It becomes the public boundary at which ordinary Python enters the stack. Surface forms such
as \py{@guarantee}, \py{assume()}, \py{check()}, and \py{hole()} are user-facing names for
moves that later become sections, obligations, and proof targets.

Zero-change verification is no longer merely a convenience mode. It becomes the inverse
boundary of the same project: instead of starting with explicit sections, the analyzer tries
to infer local sections from existing code and to transport them upward into the stack.

Theory packs are no longer a plugin mechanism with independent semantics. They are local
geometries inside the same project. Each pack says which sites matter, which solver or proof
technology can act locally there, and what sort of trust upgrade is justified when the local
problem is discharged.

The generation agent is no longer a free-floating LLM wrapper. It is a cover-growing and
cover-repairing machine. Its task is to produce local sections whose interfaces, proof
obligations, and semantics glue. Its failures are therefore not opaque generation mishaps;
they are obstructions in the stack.

\section*{Axioms of the project}
A project feels single when its commitments can be stated axiomatically. The repository is
already near that point. A reasonably faithful list of axioms is the following.
\begin{enumerate}[leftmargin=2em, itemsep=0.4em]
  \item \textbf{Locality first.} Software facts are obtained locally before they are claimed
  globally.
  \item \textbf{Types before heuristics.} Every reusable judgment should be represented, when
  possible, as a typed local section rather than a nameless heuristic.
  \item \textbf{Evidence should be explicit.} Trust levels, provenance, and residuals belong
  to the object, not merely to a log file.
  \item \textbf{Non-gluing is structured.} Bugs, contradictions, hallucinations, and failed
  synthesis attempts should be classified as organized obstructions.
  \item \textbf{Repair is refinement.} Failed globality should trigger stronger local models,
  finer covers, or better transports rather than silent fallback.
  \item \textbf{Adoption must be gradual.} The stack should admit both zero-change analysis and
  deeply annotated mixed-mode programming.
  \item \textbf{Generation belongs inside semantics.} Planning and synthesis are not beyond the
  theory; they are part of it.
\end{enumerate}

This axiomatic register is one of the main things the monograph was previously missing. Once
it is added, the later chapters stop sounding like catalog entries and start sounding like
consequences.

\section*{The repository as one mathematical object}
A final way to say the same thing is this: the repository is trying to build one mathematical
object with many interfaces. One interface is foundational and symbolic. One is operational
and proof-oriented. One is diagnostic and user-facing. One is heuristic and exploratory.
What gives the project dignity is not that each interface is already perfect. It is that each
interface can be described in terms of the same stack.

If the rest of the manuscript succeeds, the reader should reach the end feeling not that they
have toured many ambitious ideas, but that they have spent time inside one ambitious idea
that keeps differentiating itself into the forms demanded by software practice.

\chapter{Sites, Presheaves, and Locality}

Long before one speaks about hybrid dependent types, one needs a language for local
observation. The older sheaf-oriented work in the repository provides exactly that. A
program is not treated as one monolithic proof obligation. Instead it is covered by
observation sites: argument boundaries, branch guards, call results, error-sensitive
locations, output boundaries, and in the more expansive monograph, also SSA values,
wrapper boundaries, protocol or heap sites, proof sites, trace sites, and loop headers.

\begin{definition}[Program site category]
Given a program $P$, a program site category $\Site(P)$ is a category whose objects are
observation sites and whose morphisms are semantic transport maps induced by data flow,
control flow, wrapping, projection, or proof transport. A cover of a region of the
program is a family of sites whose joint information is intended to describe that region.
\end{definition}

The point of a site category is not category-theory theatrics. It is an implementation
stance. Once observations become sites, the analyzer can assign typed local sections to
those sites and restriction maps to the arrows between them. For example, an argument site
can carry a refinement about non-emptiness, a guard site can carry a branch-sensitive
fact, a call-result site can carry a transported summary from another function, and an
error site can carry a viability predicate describing when an operation is defined. The
question ``does the function type-check?'' is replaced by the more informative question
``do these local sections glue to a global contract?''

\begin{definition}[Semantic presheaf]
A semantic presheaf on a program site category is a contravariant functor
\[
  \Sem : \Site(P)^{\op} \to \mathsf{Lat}
\]
that assigns to each site a lattice of local sections and to each morphism a restriction
map transporting information backward along the corresponding semantic dependency.
\end{definition}

The repository repeatedly treats local sections as refinement-typed facts with enough
structure to be merged, compared, and explained. This is already visible in the older
paper material, where local sections are the atoms of bug finding and equivalence. It is
also visible in newer code, where bridges and checkers keep trying to harvest and move
site-local obligations.

\paragraph{Why the presheaf point of view matters.}
A presheaf gives \DepPy{} three important benefits.
\begin{enumerate}[leftmargin=*, itemsep=0.4em]
  \item It makes locality explicit. Every fact lives somewhere rather than vaguely on the
  whole program.
  \item It makes transport explicit. Facts are moved along named restrictions rather than by
  invisible global heuristics.
  \item It makes non-gluing diagnosable. Failures can be localized to overlaps and expressed
  as cocycles rather than as flat error lists.
\end{enumerate}

This is the conceptual hinge between the older sheaf-cohomology work and the newer hybrid
work. The hybrid system adds a layer dimension, but the idea of local sections and
restriction maps is already present at the foundational level.

\paragraph{Running example.}
Consider a simple normalization function.
\begin{lstlisting}
def normalize(values):
    total = sum(values)
    return [v / total for v in values]
\end{lstlisting}
A local section at the argument boundary may say ``values is a list of integers.'' A
local section at the division site may require ``$total \neq 0$.'' A local section at a call
site may transport assumptions from the caller. If those local sections do not agree on
the overlap, the failure should be represented as a mismatch in the presheaf rather than
as a mysterious one-line warning. That is the prototype of the cohomological story.

\begin{remark}
The repository sometimes writes as though the phrase ``bugs are $H^1$ classes'' is an
entire theory by itself. The more exact statement is that the site/presheaf framework
creates a mathematically principled place for gluing failures to live. The slogan is a
consequence of that framework, not a substitute for it.
\end{remark}

\chapter{The Plain Type System}

The plain type system is the foundational language of \DepPy{}. Its home is
\repo{src/deppy/types/}, and its central role is to define the semantic objects that later
passes manipulate. The file \repo{src/deppy/types/base.py} opens by describing itself as a
base hierarchy for \DepPy{}'s sheaf-descent semantic typing system. That phrase is easy to
overlook, but it says something important: the foundational system is already interpreted
semantically and geometrically.

\section*{TypeBase and the carrier language}
Every concrete plain type inherits from \py{TypeBase}. The abstract interface requires
substitution, free-variable computation, structural traversal, equality, hashing, and a
human-readable representation. Those are the operations one needs if types are to be
symbolic objects in a real formal language rather than inert runtime labels.

The concrete hierarchy covers primitive Python types, top and bottom types,
container-oriented types, callables, protocols, classes, and modules. This breadth matters
because it means the plain system aims to be a semantic typing language for native Python,
not a separate toy calculus with only theorem-prover examples.

\section*{Refinement types}
The refinement layer in \repo{src/deppy/types/refinement.py} defines the familiar shape
\[
  \{v : \tau \mid \varphi(v)\}
\]
and then invests serious engineering effort in making predicates behave like real symbolic
objects. Predicates support free-variable analysis, substitution, negation, conjunction,
disjunction, implication, comparison, and named qualifiers. This is the first place where
one sees the repository's long-term ambition: specifications are not meant to be English
comments left floating around; they are intended to be structured semantic objects.

The crucial point for later chapters is that the plain refinement layer has no trust
metadata. A refinement is simply a predicate-bearing type. There is no confidence score,
no provenance, and no solver-versus-oracle distinction internal to the object itself. This
absence is deliberate. It is one of the reasons the hybrid layer later becomes necessary.

\section*{Dependent formers}
The file \repo{src/deppy/types/dependent.py} gives the foundational system four of the
most important ingredients for the whole repository:
\begin{itemize}[leftmargin=2em]
  \item \py{PiType} for dependent functions $\Pi(x : A).B(x)$,
  \item \py{SigmaType} for dependent pairs $\Sigma(x : A).B(x)$,
  \item \py{IdentityType} for equality-like objects,
  \item indexed and quantified families for parameterized semantic structure.
\end{itemize}
These are not decorative imports from dependent type theory. They are what allow the rest
of the repository to speak concretely about APIs whose outputs depend on inputs,
witness-bearing results, transport of equalities, and theorem-carrying program fragments.

\begin{definition}[Plain dependent function]
A plain dependent function type is a triple $(x, A, B)$ represented in code as
\py{PiType(param_name, param_type, return_type)}, with substitution and free-variable
operations respecting binder hygiene.
\end{definition}

\begin{definition}[Plain dependent pair]
A plain dependent pair type is a triple $(x, A, B)$ represented as
\py{SigmaType(fst_name, fst_type, snd_type)}, where the second component type may depend on
the first component.
\end{definition}

\section*{Subtyping and obligations}
The plain system does not stop at representation. The subtype checker in
\repo{src/deppy/types/subtyping.py} produces subtype answers and verification conditions.
This is what turns the type hierarchy into an ordered refinement lattice rather than a mere
syntax tree. When later hybrid code says it values types in a refinement lattice, this is
the layer providing the foundational meaning of that phrase.

\paragraph{Why the plain system still matters after hybridization.}
It is tempting to think the hybrid system replaces the plain one. That reading is wrong.
The mixed-mode compiler still lowers user-facing syntax into plain obligations such as
refinements and Sigma-like witness forms. Lean translation still needs a foundational
type-level language. The agent and checkers still benefit from a symbolic substrate with
substitution and binder structure. The correct statement is therefore:
\begin{quote}
The plain system is the object language; the hybrid system is a richer semantics for how
objects of that language are layered, checked, and trusted.
\end{quote}

\begin{proposition}
If the hybrid system were removed but the plain type system retained, \DepPy{} would still
have a meaningful dependent-and-refinement type language. If the plain system were
removed but the hybrid packaging retained, the repository would lose the exact symbolic
objects that many higher-level workflows require.
\end{proposition}

The proof is architectural rather than formal: the code base repeatedly compiles into,
imports from, or conceptually depends on the plain hierarchy.


\part{Hybrid Evidence}
\chapter{The Hybrid Type System}

The hybrid system in \repo{src/deppy/hybrid/} is best understood as a lift of the plain
system into a layered, trust-bearing semantics. The key file is
\repo{src/deppy/hybrid/core/types.py}, whose opening documentation states the central
mathematical object directly:
\[
  \HybridTy : (\Site(P) \times \Layer)^{\op} \to \mathsf{RefinementLattice} \times \mathsf{TrustLattice}.
\]
This formula does a large amount of conceptual work. It says that a hybrid type is not
just a base type plus an extra annotation. It is a section defined over a product site
whose second coordinate ranges over the ordered set
\[
  \mathrm{INTENT} < \mathrm{STRUCTURAL} < \mathrm{SEMANTIC} < \mathrm{PROOF} < \mathrm{CODE}.
\]

\section*{Layered sections}
The hybrid core defines separate section-like structures for each layer. Intent sections
store natural-language descriptions and structured claims. Structural sections store
machine-checkable predicates, often tied to Z3 or runtime evaluation. Semantic sections
store oracle-judged predicates with prompts, rubrics, and confidence thresholds. Proof
sections store Lean-facing theorem content. Code sections store executable source and
structural program representations.

The result is a model in which one and the same semantic object can be described at
multiple levels without collapsing those levels into one another. This is the central
advantage of the hybrid formulation. It gives \DepPy{} a disciplined answer to the
problem of mixed evidence.

\begin{figure}[h]
\centering
\begin{tikzpicture}[
  layer/.style={draw=chapterblue, rounded corners=2pt, fill=chapterblue!6, minimum width=6.2cm, minimum height=0.8cm, align=center},
  score/.style={draw=accentgold, rounded corners=2pt, fill=accentgold!12, minimum width=3.4cm, minimum height=0.8cm, align=center},
  >=Latex
]
\node[layer] (intent) {INTENT \\ natural-language thesis and desired behavior};
\node[layer, below=0.45cm of intent] (structural) {STRUCTURAL \\ decidable predicates, Z3 obligations, runtime checks};
\node[layer, below=0.45cm of structural] (semantic) {SEMANTIC \\ oracle-judged rubrics, confidence, provenance};
\node[layer, below=0.45cm of semantic] (proof) {PROOF \\ Lean translations, theorem objects, residual assumptions};
\node[layer, below=0.45cm of proof] (code) {CODE \\ executable Python, extracted artifacts, assembled modules};
\node[score, right=1.7cm of semantic] (trust) {trust lattice \\ weakest-link composition};
\foreach \a/\b in {intent/structural,structural/semantic,semantic/proof,proof/code} {
  \draw[->, thick, chapterblue] (\a) -- (\b);
}
\draw[->, thick, accentgold] (semantic) -- (trust);
\draw[->, thick, accentgold] (proof.east) .. controls +(1.1,0.0) and +(0.0,-0.8) .. (trust.south);
\end{tikzpicture}
\caption{The layered object behind hybrid typing in \DepPy{}. The system keeps intent, structure, semantics, proof, and code distinct while forcing them to cohere under restriction and trust composition.}
\end{figure}

\begin{definition}[Hybrid type]
A hybrid type is a tuple of layer sections
\[
  \tau = (\tau_I, \tau_S, \tau_M, \tau_P, \tau_C)
\]
with components for intent, structural, semantic, proof, and code layers, together with
restriction maps between adjacent layers and a trust label that records the strongest
validated interpretation currently available.
\end{definition}

\section*{Dependent formers survive the lift}
The hybrid layer does not abandon the ordinary dependent type formers. Instead it defines
hybrid analogues of them. The repository names \py{HybridPiType}, \py{HybridSigmaType},
and \py{HybridRefinementType}. Their point is not to replace $\Pi$ and $\Sigma$ with new
syntax, but to reinterpret them in a setting where domain and codomain can themselves carry
structural and semantic predicates, provenance, and trust labels.

This is why the hybrid system should be read as an extension rather than a fork. A hybrid
$\Pi$-type is still fundamentally a dependent function. A hybrid $\Sigma$-type is still a
witness-bearing or decomposition-bearing pair. A hybrid refinement type still says
``values of this type satisfy this predicate,'' except that the predicate may now split into
a decidable part and a semantic part:
\[
  \{x : T \mid \varphi_d(x) \land \varphi_s(x)\}.
\]

\section*{Coherence as vanishing first cohomology}
The hybrid core says a type is coherent when first cohomology vanishes:
\[
  \Hc{1}(\Layer, \HybridTy) = 0.
\]
Operationally, this means adjacent layers agree on their overlaps. A function whose code
runs but whose semantic layer rejects its behavior is not coherent. A proof artifact that
postulates a theorem stronger than what the structural layer can justify is not coherent.
A natural-language intent that is never realized by any code or proof is not coherent.

This criterion is philosophically valuable because it makes mismatch visible without
pretending every mismatch is the same kind of failure. Some failures are structural. Some
are semantic. Some are proof-level. The language of layered cocycles keeps those failures
organized.

The hybrid geometry can also be pictured as a commutative square: restrictions in the site
dimension should commute with forgetful transport between layers. For a local chart overlap
$U \subseteq V$, one wants a square of the form
\[
\begin{tikzcd}[column sep=large, row sep=large]
\tau_S(V) \arrow[r, "\rho_{I \leftarrow S}"] \arrow[d, "\mathrm{res}_{U,V}"'] &
\tau_I(V) \arrow[d, "\mathrm{res}_{U,V}"] \\
\tau_S(U) \arrow[r, "\rho_{I \leftarrow S}"] &
\tau_I(U)
\end{tikzcd}
\]
to commute. When such squares fail to commute across a cover, the mismatch is precisely
the sort of obstruction that the hybrid cohomology language is meant to expose.

\begin{remark}
The repository sometimes introduces convenience trust enums in one file and richer trust
structures in another. A precise reading is that the hybrid package exposes a lightweight
public face and a heavier core. The paper should acknowledge this rather than flatten it.
\end{remark}

\chapter{Cohomology and Obstruction Theory}

If one chapter must bear the weight of the title, it is this one. The phrase
``cohomological typing'' should not merely signal mathematical taste. It should explain why
so many parts of the repository feel related. Cohomology is the device by which one records
that local compatibility problems are not noise but structure. In this monograph, that means
that the difference between success and failure is not exhausted by a boolean result. Success
is the existence of a global section; failure is the presence of a cocycle that refuses to be
resolved by the currently chosen cover.

\section*{From local sections to Čech complexes}
Let a cover of the relevant program object be written as \(\mathcal{U} = \{U_i\}_{i \in I}\),
and let \(\mathcal{F}\) be the presheaf of typed local evidence under discussion. Then the
first terms of the Čech complex are
\[
0 \to \Cech^0(\mathcal{U}, \mathcal{F}) \xrightarrow{\delta^0}
\Cech^1(\mathcal{U}, \mathcal{F}) \xrightarrow{\delta^1}
\Cech^2(\mathcal{U}, \mathcal{F}) \to \cdots.
\]
A $0$-cochain is a family of local sections. Its coboundary measures disagreement on pairwise
overlaps. A vanishing coboundary says that the family is compatible; a nontrivial class in
\(\Hc{1}(\mathcal{U},\mathcal{F})\) says that the local data does not glue globally.

For the software situations emphasized in this repository, one often works with covers whose
salient failures are already visible on pairwise overlaps. This is why first cohomology keeps
reappearing. The point is not that higher cohomology is philosophically impossible. The point
is that the repository's main engineering pathologies---mismatched contracts, inconsistent
local semantics, ungluable modules, competing trust claims---already have a natural home in
\(\Hc{1}\).

\section*{Global sections, obstructions, and repair}
The correct moral is therefore short:
\begin{quote}
\emph{A successful verification or synthesis episode is a vanishing first-obstruction claim
relative to a chosen cover and chosen evidence presheaf.}
\end{quote}

This sentence is powerful because it scales. For a plain typing problem, the cover may be a
family of control-flow and data-flow sites, and the presheaf may record refinement
obligations. For a hybrid coherence problem, the cover is enriched by a layer dimension and
the presheaf now records both obligations and trust-bearing witnesses. For a synthesis
problem, the cover is a family of local module obligations. For anti-hallucination, the cover
contains symbol-resolution, contract, and behavior sites. The outer mathematical sentence is
unchanged.

Repair then becomes a map from obstruction classes to refined local structure. In words,
repair is the act of changing the cover, strengthening a local predicate, refining a module
boundary, or introducing a new witness so that the old cocycle becomes a coboundary. This is
why the repository often treats diagnosis and repair as intimately connected. Once the
obstruction has been localized, one already knows where a refinement should be attempted.

\section*{Five manifestations of first cohomology}
The promise of a single project becomes most visible when the same cohomological sentence is
instantiated in several apparently distant workflows.

\paragraph{Bugs.}
In the most classical case, a bug is a non-gluing local fact. A guard site may allow states
that an arithmetic site forbids. A call site may assume a postcondition stronger than what
the callee provides. A proof site may rely on a branch invariant that the code never really
establishes. These disagreements live on overlaps and therefore naturally become $1$-cocycles.

\paragraph{Hallucinations.}
A generated artifact can be structurally well formed yet fail to inhabit the semantic layer,
or it can fail already at the structural layer because it names nonexistent symbols or
grossly misstates a signature. In either case the failure is local-to-global noncoherence.
The artifact does not inhabit the expected hybrid presheaf across the relevant cover. This is
a cohomological failure, not merely a qualitative disappointment.

\paragraph{Equivalence failure.}
Two implementations may agree at many local sites and still fail to produce a global
isomorphism. The equivalence machinery in \repo{src/deppy/equivalence/} becomes much easier
to understand once local semantic correspondences are read as candidate $0$-cochains whose
incompatibilities define a difference cocycle.

\paragraph{Generation failure.}
A module plan is a cover. Local code generation attempts to populate that cover with sections.
If interfaces, assumptions, proof goals, or exported semantics fail to match across overlaps,
one does not merely have ``bad generation.'' One has an obstruction to descent. The agent is
then in the business of refining covers and retrying gluing.

\paragraph{Trust inconsistency.}
Perhaps the most novel case is trust. One local path may deliver Z3-backed evidence, another
only oracle-backed evidence, and a third may contradict both. These are not just numbers on a
dashboard. They are typed local witnesses whose incompatibility can itself be treated as a
non-gluing phenomenon. The language of trust-valued presheaves is therefore not optional if
one wants the system to reason honestly about mixed evidence.

\begin{figure}[h]
\centering
\begin{tikzpicture}[
  box/.style={draw=chapterblue, rounded corners=2pt, fill=chapterblue!6, minimum width=3.2cm, minimum height=0.9cm, align=center},
  obs/.style={draw=accentgold, rounded corners=2pt, fill=accentgold!12, minimum width=3.3cm, minimum height=0.9cm, align=center},
  >=Latex
]
\node[box] (local) {local sections \ typed evidence on sites};
\node[box, right=1.6cm of local] (overlap) {overlap maps \ pairwise compatibility};
\node[obs, right=1.6cm of overlap] (h1) {$\Hc{1}$ class \ unresolved obstruction};
\node[obs, below=1.0cm of local] (bugs) {bugs};
\node[obs, below=1.0cm of overlap] (hall) {hallucinations / equivalence failures};
\node[obs, below=1.0cm of h1] (gen) {generation / trust inconsistencies};
\draw[->, thick, chapterblue] (local) -- (overlap);
\draw[->, thick, chapterblue] (overlap) -- (h1);
\draw[->, thick, accentgold] (h1) -- (bugs);
\draw[->, thick, accentgold] (h1) -- (hall);
\draw[->, thick, accentgold] (h1) -- (gen);
\end{tikzpicture}
\caption{One cohomological sentence, many manifestations. The repository's main failure modes are organized by the same passage from local sections to overlap incompatibilities to first-obstruction classes.}
\end{figure}

\section*{Obstruction mass and the quality of an elaboration}
The repository's newer scalar-valued ideation layer is also easier to justify in this
language. A candidate elaboration should not be judged only by thematic elegance or by an
LLM's enthusiasm. It should also be penalized when it appears likely to generate unresolved
obstruction mass. In the current code, this appears as a penalty term related to the
observed or estimated \(H^1\)-dimension of the ideation result. That move is mathematically
more principled than it first appears. It says that elaborations carrying unresolved gluing
risk should be scored lower even before full generation begins.

In other words, cohomology does double duty. It diagnoses explicit failures after local
claims have been made, and it also helps rank prospective elaborations before the full cost
of construction is paid. This is one of the clearest signs that the project's geometry is
not ornamental. It actively shapes the agent's search.

\section*{Why first cohomology is not only about failure}
There is a temptation to speak as though cohomology entered the story only when something has
gone wrong. That reading is too narrow. The vanishing of first cohomology is not merely the
absence of drama. It is the positive condition that local evidence has achieved global form.
A gluable family of sections is itself a mathematical achievement. It is a certificate that
multiple local accounts of a program agree. For a system like \DepPy{}, which deliberately
keeps many kinds of evidence visible instead of collapsing them, such certificates are the
closest thing to semantic harmony that one can ask for.

\section*{Repair as a research methodology}
This suggests a wider philosophical point. The repository is not only trying to type and
verify programs. It is trying to organize research methodology around the cycle
\[
  \text{observe locally} \to \text{compute incompatibility} \to \text{refine locally} \to \text{glue globally}.
\]
That cycle is visible in program analysis, in theory-pack extension, in ideation, and in
agent-driven synthesis. Once stated this way, it becomes plausible to claim that
``cohomological typing for absolutely everything'' is a methodology and not merely a slogan.

\begin{proposition}
If a subsystem of software engineering can be represented by a cover of local sites, a notion
of typed local evidence, and a repair loop that attempts to turn overlap failures into
compatibilities, then that subsystem is in scope for the repository's cohomological typing
program.
\end{proposition}

The proof is not a theorem internal to one file; it is the architecture of the repository
itself.

\chapter{Trust, Oracles, and Proofs}

The hybrid system's most distinctive engineering move is to treat semantic judgment as a
controlled effect rather than as an invisible side channel. This is expressed by the
oracle monad
\[
  \TO(A) = A \times \mathsf{TrustLevel} \times \mathsf{Confidence} \times \mathsf{Provenance}.
\]
A semantic check is therefore not merely a boolean returned by an LLM. It is a value with
a trace: what model or process judged it, how confident that judgment is, whether it was
cached, and which trust level should be attached to the result.

\section*{The trust lattice}
The core trust ordering used by the hybrid types can be written as
\[
  \mathrm{CONTRADICTED} < \mathrm{UNTRUSTED} < \mathrm{RUNTIME\_CHECKED} < \mathrm{LLM\_JUDGED} < \mathrm{Z3\_PROVEN} < \mathrm{LEAN\_VERIFIED}.
\]
Some convenience exports use slightly different public names, but the semantic pattern is
the same: trust is ordered, and the order is used algebraically.

The meet operation records weakest-link composition,
\[
  a \sqcap b = \min(a,b),
\]
while the join operation records strongest evidence,
\[
  a \sqcup b = \max(a,b).
\]
These operations matter because they make compositionality honest. If one stage is only
LLM-judged while a downstream stage is solver-proven, the composition should not suddenly
pretend to be Lean-verified. The weakest link constrains the trust of the whole.

\section*{Structural versus semantic predicates}
The oracle monad would be much less useful if every predicate lived inside it. The
repository instead distinguishes structural predicates from semantic ones. Structural
predicates are intended to be decidable or at least deterministically checkable. Semantic
predicates encode claims whose meaningful evaluation may require richer interpretation,
rubrics, or domain knowledge.

This split is the deepest reason the hybrid system can talk about anti-hallucination,
solver discharge, and proofs in one breath without conflating them. Structural checks may
be discharged by Z3 or deterministic evaluation. Semantic checks enter the oracle monad
and produce confidence-bearing evidence. Proof translation may then reflect some of these
facts into Lean, leaving others as assumptions or residual obligations.

\begin{definition}[Hybrid refinement]
A hybrid refinement type is a base type $T$ equipped with a structural predicate
$\varphi_d$ and an optional semantic predicate $\varphi_s$, written
\[
  \{x : T \mid \varphi_d(x) \land \varphi_s(x)\}.
\]
The intent is that $\varphi_d$ is decidable while $\varphi_s$ is judged via the oracle
pipeline and therefore carries trust and confidence metadata.
\end{definition}

\section*{Lean translation and residual honesty}
The repository's Lean translation and discharge cascade are significant because they
preserve the distinction between theorem-backed and oracle-backed content. The discharge
cascade typically moves through trivial discharge, Z3, oracle-assisted proof search,
Lean checking, and residual fallback. This is an unusually honest design. It avoids the
common temptation to treat all successful pipeline stages as equally formal.

A semantic predicate may be accepted at the semantic layer and even exported into a proof
pipeline, but if it ultimately becomes an axiom or a residual assumption that boundary
should remain visible. One of the strongest lessons from the repository is therefore that
proof engineering and trust engineering must coexist. A system that hides its residuals is
less trustworthy than one that openly classifies them.

\begin{proposition}
Any semantics-preserving translation from hybrid types to Lean must either prove,
postulate, or residualize every semantic predicate that lacks a purely structural proof.
\end{proposition}

This proposition is conceptually obvious, but it is exactly the sort of obvious fact that
software tools often ignore in practice. \DepPy{} does not ignore it; the hybrid design is
built around it.

\chapter{Obstructions, Equivalence, and Hallucination}

One of the repository's strongest unifying ideas is that several seemingly different
problems become gluing problems once the right local sections are chosen.

\section*{Bugs as obstruction classes}
In the sheaf-cohomological line of work, a bug is represented as a non-gluing local fact.
A local section at one site may allow a state that another site's viability condition
forbids. The disagreement lives on an overlap and can be collected into a $1$-cocycle.
When that cocycle is not a coboundary, it determines a nonzero class in first cohomology.
This is the precise sense in which the repository says ``bugs are $H^1$ classes.''

\begin{example}[Normalization bug]
Suppose an argument site allows an empty list while a downstream division site requires a
nonzero sum. The call and error sites disagree on their overlap. The mismatch is a
$1$-cocycle; if no local strengthening already present in the cover resolves it, the class
is nonzero and the bug is genuinely obstructive.
\end{example}

The slogan becomes especially valuable once it is quantified. The repository's older paper
material argues that the rank of $\Hc{1}$ over $\mathbb{F}_2$ corresponds to the minimum
number of independent fixes required to resolve the current family of gluing failures.
Whether or not every implementation path computes that exact number in the same way, the
architectural point is powerful: non-gluing is not just yes/no, it has structure and can
be decomposed.

\section*{Equivalence as gluing of local isomorphisms}
The equivalence machinery in \repo{src/deppy/equivalence/} makes the same move for program
comparison. Two programs are not compared by one magical global theorem from nowhere.
Instead, one constructs local correspondences, local semantic differences, and overlap
conditions. If the local isomorphisms glue, one obtains a global equivalence verdict. If
they do not, one obtains a localized reason for failure.

This reuse of the same geometry is one of the best reasons to take the repository's
mathematical self-description seriously. It would be much easier to implement bug finding
and equivalence with unrelated vocabularies. The decision not to do so shows that the
local-to-global picture is doing real architectural work.

\section*{Hallucination as type inhabitation failure}
Anti-hallucination is where the hybrid split between structural and semantic predicates
becomes especially visible. A fabricated API is a structural failure: the relevant symbol,
signature, or object does not exist. A semantic drift is a meaning failure: the code might
be structurally well-formed while still failing the intended rubric or behavior. A
contradiction between multiple checks is an inconsistency failure.

This is a very strong and very repository-specific claim. It says that hallucination should
not be treated as mystical stochastic misbehavior. It should be treated as typed failure.
The artifact has a required hybrid type. Structural checks attempt to inhabit the
structural component. Semantic checks attempt to inhabit the semantic component. Proof and
code layers contribute their own evidence. If the object fails to inhabit the required
space, the failure can be named precisely.

\begin{definition}[Hallucination as inhabitation failure]
Given an expected hybrid type $T$, a generated artifact $a$ hallucinates when $a$ fails to
inhabit $T$ at some required layer. Structural hallucinations fail structural inhabitation;
semantic hallucinations fail semantic inhabitation; contradictory hallucinations produce
incompatible evidence across layers.
\end{definition}

This is one of the clearest examples of the paper's title. Once local-to-global typed
structure is the common language, hallucination, equivalence, and bug finding are not
unrelated downstream products. They are different ways for a typed section to succeed or
fail to glue.

\chapter{Theory Packs and Local Solver Geometries}

The repository does not attempt to solve every semantic question with one universal logic.
Instead it provides domain-specific theory packs and theory-library machinery. This is
another place where the cohomological typing viewpoint clarifies what the implementation is
doing.

A theory pack is best understood as a local geometry of solvable questions. Tensor shape,
nullability, boolean guards, witness transport, exceptions, heaps, protocols, loop
invariants, and similar concerns each induce their own families of local sections and
viability checks. The pack decides which sites are relevant, what local solver to invoke,
which restrictions are legitimate, and how trust should be upgraded when a local discharge
succeeds.

\section*{Why packs fit the site-based viewpoint}
The older monograph on sheaf-descent semantic typing already pushes hard on the idea of
site families and theory families. The hybrid and theory-pack code carries the same spirit
forward. Rather than pretending there is one maximally expressive theorem prover inside the
project, \DepPy{} distributes reasoning across packs. Each pack is responsible for a local
stalk theory.

That phrasing is not decorative. A stalk theory answers local questions about a site. A
shape theory can answer tensor-shape compatibility questions. A guard theory can answer
branch-sensitive path-viability questions. A protocol or heap theory can answer ownership
and field-availability questions. A witness transport theory can answer equality and
transport questions. The global analyzer then composes these local answers.

\section*{Trust upgrades through local discharge}
Theory packs also give a concrete answer to the question ``where do trust upgrades come
from?'' They come from successful local discharge in a domain where the relevant solver or
proof method is appropriate. A structural fact that begins as merely asserted may rise to
runtime-checked or Z3-proven status after a local pack solves it. The hybrid trust lattice
therefore does not float abstractly above the repository; it is fed by actual local
reasoning modules.

\begin{example}[Tensor theory]
A tensor-shape pack may certify that two tensors can be broadcast together, or that a
reshape preserves total size. Those are local shape facts, but once transported across
call sites and composed with code and proof layers they become part of a global contract
for a function. The same story generalizes to device compatibility, indexing safety, and
order-sensitive tensor operations.
\end{example}

\section*{Absolutizing the title carefully}
If the phrase ``absolutely everything'' is to remain honest, theory packs are part of the
answer. The repository cannot yet reason deeply about every possible Python phenomenon.
But it can repeatedly grow new local geometries and plug them into a common architecture.
The title should therefore be read as a claim about the architecture's extensibility: any
domain that can be given local sections, restrictions, and gluing conditions is in scope
for cohomological typing.


\part{Elaboration and Synthesis}
\chapter{Type Expansion, Ideation, and Agents}

The repository's prompt-driven agent is not an accidental extra on top of the typing
system. It is one of the clearest demonstrations that the same semantic stance can govern
planning and generation as well as checking.

\section*{Type expansion}
The type-expansion modules in the hybrid package take a short natural-language description
and elaborate it into a much richer plan: modules, covers, dependencies, obligations,
interfaces, and eventually code-generation structure. This process is easiest to interpret
as growth in a presheaf. The intent layer starts with sparse, natural-language local
sections. Structural, semantic, proof, and code layers gradually acquire more content as
the plan becomes more explicit.

The important conceptual shift is that generation is not described as free-form text
completion followed by a verifier. It is described as typed elaboration. The prompt gives
an intent section. Expansion grows a cover. Local generation produces candidate sections.
Verification checks whether they glue. Repair loops try to resolve obstructions. Assembly
produces a global artifact.

\section*{Ideation as typing}
The ideation machinery is even more unusual. The repository includes a large mathematical
ontology and domain-selection logic intended to enrich prompts by analogy to concepts from
many mathematical areas. The right way to phrase this in the paper is not that the system
has a magical creativity engine. It is that it tries to treat ideation itself as a typed
search over richer local descriptions.

If a prompt about trading is enriched by analogy to algebraic geometry, order theory,
signal processing, or topology, the resulting design search space is larger but also more
structured. Ideation becomes a pre-typing phase that improves the quality of the later
cover and the plausibility of the resulting local sections.

\section*{Scalar-valued typing}
The ideation pipeline can be sharpened still further by introducing what the repository now
calls \emph{scalar-valued typing}. The idea is simple: before committing to a full
elaboration, assign a scalar to the candidate elaboration that measures how strongly the
current domain evidence supports it. In implementation terms, \repo{agent/cli.py} already
tracks an LLM belief score, a mathematically-derived score based on validated analogies and
certifiable structure, and per-area strengths. Scalar-valued typing packages these into a
single typed quantity attached to the elaboration.

If $E$ is a type elaboration and $d$ is a target domain, then the scalar-valued typing
score can be written schematically as
\[
  \Score(E,d) =
  \alpha \, \Score_{\mathrm{LLM}}(E,d)
  + (1-\alpha)\, \Score_{\mathrm{Math}}(E,d)
  + \varepsilon \, \max_i \Score_i(E,d),
\]
with \(\alpha = 0.55\) and \(\varepsilon = 0.05\) in the current implementation. The
first term measures how strongly the language model believes the idea explains the problem
in the chosen domain. The second term measures structural richness: validated analogies,
deep connections, theorem schemas, certification targets, and penalties for unresolved
\(H^1\)-mass. The final term lets especially strong local mathematical areas gently boost a
globally promising elaboration.

This is a useful addition because it makes ideation behave less like unranked brainstorming
and more like typed search. Scalar-valued typing does not replace ordinary type checking.
Instead it orders candidate elaborations before the expensive downstream steps of full
module expansion, proof planning, and code generation. The score is therefore neither a
proof nor a mere UX flourish; it is a typed prioritization signal.

\begin{figure}[h]
\centering
\begin{tikzpicture}[
  node distance=1.15cm,
  box/.style={draw=chapterblue, rounded corners=2pt, fill=chapterblue!6, minimum width=3.1cm, minimum height=0.9cm, align=center},
  sum/.style={draw=accentgold, rounded corners=2pt, fill=accentgold!12, minimum width=3.5cm, minimum height=0.95cm, align=center},
  >=Latex
]
\node[box] (intent) {prompt / intent section};
\node[box, right=of intent] (llm) {LLM belief \\ \(\Score_{\mathrm{LLM}}\)};
\node[box, right=of llm] (math) {math-model evidence \\ \(\Score_{\mathrm{Math}}\)};
\node[box, below=of llm] (areas) {area strengths \\ \(\max_i \Score_i\)};
\node[sum, below=of math] (overall) {scalar-valued typing \\ \(\Score(E,d)\)};
\draw[->, thick, chapterblue] (intent) -- (llm);
\draw[->, thick, chapterblue] (intent.east) .. controls +(1.5,-0.4) and +(-1.4,0.6) .. (math.west);
\draw[->, thick, chapterblue] (llm) -- (overall);
\draw[->, thick, chapterblue] (math) -- (overall);
\draw[->, thick, chapterblue] (areas) -- (overall);
\end{tikzpicture}
\caption{Scalar-valued typing adds a numeric layer over ideation: type elaborations are scored by domain fit before full expansion.}
\end{figure}

\section*{The orchestration agent}
The agent orchestrator in \repo{agent/orchestrator.py} embodies a multi-stage state
machine: parsing, expanding, planning, generating, verifying, and assembling. The code is
therefore easier to understand after one adopts the hybrid type picture. Each stage is not
simply passing strings to the next. It is refining what kind of section exists, which
layer it inhabits, and which trust-bearing obligations must be checked before the next
stage is allowed to treat the result as global enough.

\begin{algorithm}[h]
\caption{Cohomological generation loop}
\begin{algorithmic}[1]
\State Resolve user prompt into intent sections.
\State Expand intent into a typed cover of local module obligations.
\For{each local module site}
  \State Synthesize candidate code and local contracts.
  \State Check structural, semantic, and proof obligations.
  \If{an obstruction remains}
    \State refine the local section and retry
  \EndIf
\EndFor
\State Attempt to glue the verified local sections into a project.
\State Emit residual obstructions if gluing fails.
\end{algorithmic}
\end{algorithm}

The importance of this loop is that it places generation inside the same semantics as
checking. A paper that really integrates the repository should emphasize this strongly. In
the current repository story, scalar-valued typing sits just before the expensive center of
this loop, biasing the search toward elaborations whose local mathematical evidence most
strongly supports the target domain.

\chapter{Native Python, Zero-Change Checking, and Mixed Mode}

One of the repository's recurring practical commitments is that users should not be forced
into a theorem-prover syntax cliff before they receive useful results. This commitment has
two implementation faces: zero-change checking for existing code and mixed-mode syntax for
new code.

\section*{Zero-change verification}
The package-level quick APIs and zero-change tools are designed so ordinary Python can be
analyzed without being rewritten first. The repository's exports around
\py{hybrid_check()}, \py{hybrid_analyze()}, and implicit presheaf extraction form a gentle
entry point into the larger architecture. This is important because it shows that the
hybrid system is not reserved for carefully annotated greenfield projects.

The right way to read zero-change verification is as inference of local sections from
existing program structure. The analyzer extracts sites from source, docstrings, control
flow, and library use, then runs structural and semantic passes over those inferred sites.
Trust labels remain visible, so the user can distinguish inferred, solver-backed, and
residual facts.

\section*{Mixed-mode syntax}
The mixed-mode syntax in \repo{src/deppy/hybrid/mixed_mode/syntax.py} is the converse
adoption path. Instead of starting from bare existing code, one starts from ordinary Python
augmented with controlled surface forms such as \py{hole()}, \py{spec()},
\py{@guarantee}, \py{assume()}, and \py{check()}. The crucial invariant named in the file
is graceful degradation: before compilation the forms either pass through, return safe
defaults, or raise a controlled placeholder exception. This allows editors, tests, and
ordinary Python execution to remain useful.

\begin{lstlisting}
from deppy.hybrid import guarantee, assume, check, hole

@guarantee("returns a sorted list with no duplicates")
def unique_sorted(xs: list[int]) -> list[int]:
    assume("xs is finite")
    ys = hole("remove duplicates from xs")
    check("ys preserves the elements of xs")
    return sorted(ys)
\end{lstlisting}

The compiler path for these forms reveals the deep relation between the plain and hybrid
systems. Surface forms compile into plain obligations such as refinements and Sigma-like
witness objects, then re-enter the hybrid pipeline as trust-bearing contracts and proof
obligations. This makes mixed mode the public handshake between the two type systems.

\section*{Why this matters for the paper's title}
If cohomological typing is really ``for absolutely everything,'' it must include both
legacy programs and new proof-oriented code. Zero-change checking handles the first case;
mixed mode handles the second. Together they show that the architecture is not only a
research kernel but also a migration strategy.

\chapter{Examples, Trading, and Repository Integration}

A monograph about repository integration should not stay at the level of definitions. It
should show how those definitions illuminate the code that users actually see.

\section*{Examples as semantic stress tests}
The examples shipped in the repository exercise several different faces of the system:
proof-oriented examples, tensor and shape examples, mixed-mode examples, zero-change
checking examples, and the trading examples that drive the prompt-based generation story.
These are not merely demos. They are semantic stress tests for the thesis that one local
to-global architecture can support multiple workflows.

\section*{Why the trading system matters}
The trading story is useful because it forces the repository to coordinate many ideas at
once. There is natural-language intent about discovering or constructing a strategy. There
are structural constraints such as position sizing, leverage bounds, exposure rules, and
interface boundaries between market data, strategy kernels, risk engines, execution, and
portfolio state. There are semantic judgments about whether a strategy meaningfully realizes
a stated thesis. There are proof ambitions around risk-critical invariants. And there is
actual code generation and assembly.

The hybrid cover-oriented generation pipeline is therefore not theoretical ornament in the
trading example. It is exactly the mechanism needed to keep multiple modules and multiple
kinds of evidence from drifting apart. A trading project becomes a cover of local sections:
market-data sections, signal sections, risk sections, execution sections, portfolio
sections, compliance sections, and final assembly. The question ``did generation succeed?''
becomes a gluing question.

\begin{figure}[h]
\centering
\begin{tikzpicture}[
  node distance=0.95cm,
  site/.style={draw=chapterblue, rounded corners=2pt, fill=chapterblue!6, minimum width=2.35cm, minimum height=0.8cm, align=center},
  >=Latex
]
\node[site] (market) {market\_data};
\node[site, right=of market] (signal) {signal\_geometry};
\node[site, right=of signal] (kernel) {strategy\_kernel};
\node[site, below=of signal] (risk) {risk\_engine};
\node[site, right=of risk] (exec) {execution};
\node[site, right=of exec] (portfolio) {portfolio};
\node[site, below=of exec] (compliance) {compliance};
\node[site, below=of kernel] (main) {main / global section};
\draw[->, thick, chapterblue] (market) -- (signal);
\draw[->, thick, chapterblue] (signal) -- (kernel);
\draw[->, thick, chapterblue] (kernel) -- (risk);
\draw[->, thick, chapterblue] (risk) -- (exec);
\draw[->, thick, chapterblue] (exec) -- (portfolio);
\draw[->, thick, chapterblue] (portfolio) -- (main);
\draw[->, thick, chapterblue] (compliance) -- (main);
\draw[->, dashed, accentgold] (risk) -- (compliance);
\draw[->, dashed, accentgold] (kernel) -- (main);
\end{tikzpicture}
\caption{A repository-native trading cover. Solid arrows follow the execution and state flow; dashed arrows indicate proof and compliance constraints that must still glue into the final global section.}
\end{figure}

The newer scalar-valued typing layer fits naturally into this trading picture. Different
mathematical lenses---for example stochastic calculus, financial mathematics, or control
theory---can assign different scalar weights to the same type elaboration. Those weights do
not prove the system correct, but they help the ideation stage decide which elaborations
deserve expansion into the expensive multi-module cover.

\section*{Repository-wide integration map}
The repository can be read through the following semantic map.
\begin{itemize}[leftmargin=2em, itemsep=0.5em]
  \item \repo{src/deppy/types/}: the foundational object language.
  \item \repo{src/deppy/hybrid/core/}: layered type semantics, trust, checker, and geometry.
  \item \repo{src/deppy/hybrid/mixed_mode/}: user-facing bridge from Python syntax to obligations.
  \item \repo{src/deppy/hybrid/discharge/} and \repo{lean_translation/}: proof-facing boundary.
  \item \repo{src/deppy/hybrid/type_expansion/}: planning and architectural elaboration.
  \item \repo{src/deppy/hybrid/anti_hallucination/}: structural and semantic artifact validation.
  \item \repo{src/deppy/equivalence/}: local-to-global comparison of programs and predicates.
  \item \repo{src/deppy/render/}: turning typed and trust-bearing results into readable diagnostics.
  \item \repo{src/deppy/theory_packs/} and theory libraries: modular local solver geometries.
  \item \repo{agent/}: the orchestration loop that tries to grow a global section from prompt-local material.
\end{itemize}

Once this map is visible, the repository stops looking like a pile of unrelated prototypes.
It looks like a large, uneven, but conceptually coherent attempt to turn local semantics
into a universal organizing principle.

\chapter{Functorial Semantics: Rendering, Transport, and Synthesis}

A unified project should not reserve its mathematics for only one corner of the system. If
local-to-global semantics is truly the repository's common language, then even apparently
secondary concerns such as rendering, transport, and synthesis should become legible in that
language. This chapter argues that they do. The right conceptual move is to treat them as
functorial constructions on typed presheaves rather than as utility layers sitting outside
the theory.

\section*{Rendering as a natural transformation}
Suppose \(\mathcal{F}\) is a presheaf of typed local evidence. A renderer does not discover a
new semantic fact ex nihilo. It reorganizes the existing typed content into another target
language: English explanation, Lean code, diagrams, tables, or machine-readable reports.
That suggests the right abstraction. Rendering is a natural transformation from a presheaf of
typed sections into a presheaf of representations.

In the repository, this idea appears in multiple places. The web-facing explanations in
\repo{web/} are not external to the type story; they are human-readable images of it. The
Lean translation pipeline in \repo{src/deppy/hybrid/lean_translation/} is not just export; it
is a semantics-preserving map from one typed universe into another. The report and rendering
machinery in \repo{src/deppy/render/} likewise maps trust-bearing objects into surfaces that a
user or auditor can inspect.

What matters is that a good renderer respects restriction maps. If one moves from a whole
function to a branch site or from a module to an overlap interface, the rendered object
should continue to agree with the underlying semantics. That is exactly the sort of square
one expects of a natural transformation.

\[
\begin{tikzcd}[column sep=large, row sep=large]
\mathcal{F}(V) \arrow[r, "R_V"] \arrow[d, "\mathrm{res}_{U,V}"'] & \mathcal{R}(V) \arrow[d, "\mathrm{res}_{U,V}"] \\
\mathcal{F}(U) \arrow[r, "R_U"] & \mathcal{R}(U)
\end{tikzcd}
\]

Here \(\mathcal{R}\) might be a Lean-syntax presheaf, a textual-report presheaf, or a diagram
presheaf. The point is the same in every case: rendering should commute with locality.

\section*{Transport as restriction and extension}
Interprocedural analysis and contract propagation are often discussed as if they were
orthogonal to typing. In this project they are not. They are the mechanics by which typed
information moves along a program's semantic edges. One may call that restriction,
extension, transport, or pullback depending on which aspect one wants to emphasize, but it
is still the same local-to-global geometry.

When a caller invokes a callee, one restricts attention from the caller's current state to
the callee's input interface. When the callee returns, one transports postconditions back
into the caller's context. When a wrapper preserves an invariant, one is witnessing a
particular compatibility between the wrapped and unwrapped sites. When a zero-change analyzer
infers a likely precondition from existing source, it is harvesting transportable local
facts from ordinary Python and lifting them into typed form.

This is why the repository's adoption story remains coherent. The same semantics that govern
explicitly annotated code can also govern inferred local structure. The difference is not one
of basic ontology but of how much evidence is given up front.

\section*{Synthesis as cover elaboration}
Perhaps the strongest test of unity is code synthesis. If synthesis sits outside the theory,
then the project fractures: mathematics on one side, generation on the other. The repository
is most interesting precisely because it refuses that fracture. The generation agent is best
understood as an elaborator of covers.

A user prompt does not immediately produce a monolithic program. It produces intent-level
material from which a cover is elaborated: modules, interfaces, invariants, proof targets,
and gluing obligations. Local generation then attempts to produce sections on each chart.
Checking passes test whether those sections agree on overlaps. Residual failures trigger
cover refinement or local repair. Only at the end does one speak about a global artifact.

This account is far superior to saying that the repository simply ``uses an LLM and then
verifies the output.'' That description misses the architecture. The architecture says that
synthesis itself is a typed local-to-global process.

\begin{figure}[h]
\centering
\begin{tikzpicture}[
  node distance=1.0cm,
  box/.style={draw=chapterblue, rounded corners=2pt, fill=chapterblue!6, minimum width=2.8cm, minimum height=0.85cm, align=center},
  obs/.style={draw=accentgold, rounded corners=2pt, fill=accentgold!12, minimum width=3.0cm, minimum height=0.85cm, align=center},
  >=Latex
]
\node[box] (prompt) {prompt};
\node[box, right=of prompt] (cover) {cover elaboration};
\node[box, right=of cover] (sections) {local sections};
\node[obs, right=of sections] (glue) {global artifact or obstruction};
\node[obs, below=of sections] (repair) {cover refinement / local repair};
\draw[->, thick, chapterblue] (prompt) -- (cover);
\draw[->, thick, chapterblue] (cover) -- (sections);
\draw[->, thick, chapterblue] (sections) -- (glue);
\draw[->, thick, accentgold] (glue) -- node[right] {if non-gluing} (repair);
\draw[->, thick, accentgold] (repair.west) .. controls +(-3.0,0.0) and +(0.0,-0.8) .. (cover.south);
\end{tikzpicture}
\caption{Synthesis as cover elaboration. The agent is most coherently understood as a machine for growing, testing, and refining covers until local sections glue.}
\end{figure}

\section*{Lean translation as a deep embedding}
The Lean pipeline deserves special mention because it dramatizes the project's ambition. A
lightweight renderer merely paraphrases typed content. Lean translation aims higher. It tries
to embed a portion of the repository's typed world into a theorem-proving environment whose
kernel can then check part of the resulting content. This is not a casual format conversion.
It is a deep embedding of selected structural and proof-bearing aspects of the stack.

That also explains why residual honesty matters so much. A functor into Lean cannot simply
pretend that every semantic predicate has become a theorem. Some content becomes theorem
statements. Some becomes lemmas with external evidence. Some becomes axioms. Some remains as
residual assumptions. A faithful deep embedding preserves these distinctions rather than
obscuring them.

\section*{Why this strengthens the sense of one project}
The unifying power of the present viewpoint is easiest to appreciate negatively. Without it,
rendering looks like documentation, transport looks like static-analysis plumbing, and
synthesis looks like an attached generation system. With it, they become further instances of
a single mathematical habit: typed local material should move coherently between sites,
between layers, and between external representations.

That is exactly the kind of broad conceptual reuse a monograph needs if it is to feel like a
single project rather than a collection of adjacent papers. The project is not merely about
which types exist. It is about how typed local evidence circulates, is rendered, is upgraded,
is repaired, and is finally assembled.

\section*{The repository as a functorial workshop}
A final image may help. The repository is not only a checker and not only a generator. It is
a workshop of functors and natural transformations built around typed presheaves. Some maps
go downward into code. Some go upward into proofs. Some go sideways into reports, diagrams,
and English explanations. Some shuttle information across call graphs and module boundaries.
What makes these maps members of one family is that they are all expected to respect the same
local-to-global semantics.

Once that expectation is made explicit, the repository reads more like a foundational
research program. It is building not just tools, but a universe in which software artifacts,
proof artifacts, semantic judgments, and generated designs can all be transported through a
single coherent language.

\chapter{Implementation Program}

This chapter turns the preceding conceptual story into a concrete implementation program
for the repository as it exists now.

\section*{Principle 1: Keep the plain system as the stable substrate}
The foundational hierarchy in \repo{src/deppy/types/} should remain the canonical place for
symbolic type-level structure: substitution, free-variable logic, refinements, dependent
formers, indexed families, and subtyping. Hybrid features should continue to depend on
that language rather than re-implementing its semantics in parallel.

\section*{Principle 2: Concentrate trust logic in a small number of core modules}
The current repository exposes both lightweight and heavyweight trust vocabularies. This is
not fatal, but it suggests an engineering direction: preserve public convenience exports,
while concentrating the real semantics of trust, provenance, and evidence composition in a
small number of clearly canonical modules. That will make papers, dashboards, and checkers
less likely to drift apart.

\section*{Principle 3: Treat covers as first-class products}
The older monographs already emphasize cover synthesis and site families. The generation
agent and hybrid pipeline should continue to move in that direction. Rather than treating
module plans or verification stages as opaque steps, they should make covers, overlap
obligations, and gluing summaries first-class artifacts. This makes the local-to-global
story inspectable by users and not only by implementers.

\section*{Principle 4: Make residuals visible}
The paper has repeatedly insisted on epistemic honesty, and the implementation should do
the same. Residual obligations, assumption boundaries, and heuristic semantic judgments
should remain visible in manifests, renderers, and reports. The trust lattice is useful
precisely because it avoids silent collapse.

\section*{Principle 5: Let theory packs grow the scope}
The phrase ``absolutely everything'' should be realized by extensibility, not by claiming
that every domain is already solved. New theory packs, domain bridges, and site families
should remain the preferred path for broadening \DepPy{}'s scope. The repository already
has the architectural shape required for this growth.

\begin{remark}
The implementation program described here is conservative in one sense and radical in
another. It is conservative because it preserves the repository's existing semantic core.
It is radical because it proposes to explain more and more of the repository through a
single local-to-global language rather than through disconnected subsystem-specific
vocabularies.
\end{remark}

\chapter{Conclusion}

The repository now contains enough material to support a genuine monograph-level thesis.
The thesis is not merely that sheaf theory is a pretty metaphor for program analysis, nor
merely that hybrid dependent types are a convenient wrapper around solver checks and LLM
judgments. The deeper thesis is that local-to-global semantics can be a universal
organizing principle across a surprising amount of software engineering work.

The plain type system already provides a strong foundational language of exact semantic
objects. The hybrid system already provides a layered semantics of trust-bearing evidence.
The geometry modules already provide a language of gluing and obstruction. The equivalence
modules already show how local isomorphism can become a global verdict. The anti-
hallucination modules already show how artifact quality can be recast as inhabitation.
Type expansion and the agent already show that planning and generation can be read as typed
elaboration. Zero-change checking and mixed mode already show that the architecture can be
adopted both from legacy Python and from proof-oriented new code.

What the revised manuscript has tried to add is not merely more detail but more inevitability.
It has tried to show that these lines of work do not just coexist; they compel one another.
Once one decides that software facts should be represented locally, that evidence should be
typed, that differences in evidence quality should remain explicit, and that non-gluing
should be treated as structure rather than noise, then a remarkable amount follows. One is
led toward refinement lattices, toward hybrid witnesses, toward trust-sensitive semantics,
toward cohomological diagnosis, toward theory packs, toward zero-change extraction, and
toward synthesis as elaboration. The unity of the project lies not in a branding exercise,
but in this chain of implication.

There is still real work ahead. The repository's different lines of development need to be
made even more consistent. Trust vocabularies need consolidation. The bridge from intent to
formalization remains heuristic and should be described as such. Covers, site families, and
residuals should be made still more explicit in user-facing outputs. But those are exactly
the kinds of tasks one expects in a large research prototype converging on a unifying
architecture.

This is also why the phrase ``absolutely everything'' should be read with discipline rather
than skepticism. It does not mean that every domain has already been mastered. It means that
the project has identified an extensible semantic shape into which new domains can be fit.
Whenever one can identify local sites, typed local evidence, transport maps, and a sensible
notion of gluing, one has a candidate entrance into the same universe. The right way to
extend the system is therefore not to bolt on one more isolated feature, but to ask what the
local sections are, where their overlaps live, what trust structure they require, and how
their obstructions should be repaired.

If the phrase ``cohomological typing for absolutely everything'' sounded excessive at the
start, the point of this monograph is to make it sound less like excess and more like a
research program. The program is this: whenever a subsystem of software work can be
understood as local evidence that should glue globally, \DepPy{} asks whether typing,
trust, and cohomology can be used to express that subsystem in one semantic language. Much
of the repository is already an existence proof that the answer is often yes.

The deepest hope behind the project is therefore not just that particular theorems will be
proved or that particular agents will generate better code. It is that software engineering
can acquire a more stable conceptual center. In this monograph, that center is the claim
that types, evidence, proofs, semantics, and generation are not unrelated vocabularies but
progressive refinements of one local-to-global discipline. If the reader leaves with that
impression, the document will have done what a foundational document should do: make later
work feel less optional and more like the next chapter of the same idea.


\appendix
\chapter{Appendix: Candidate Definitions and Results}

This appendix collects paper-ready definitions and results that fit the repository's
current architecture.

\begin{definition}[Hybrid global section]
Let $P$ be a program and let $\Layer$ be the ordered set
$\mathrm{INTENT} < \mathrm{STRUCTURAL} < \mathrm{SEMANTIC} < \mathrm{PROOF} < \mathrm{CODE}$.
A hybrid global section for $P$ is a section of the presheaf
\[
  \HybridTy : (\Site(P) \times \Layer)^{\op} \to \mathsf{RefinementLattice} \times \mathsf{TrustLattice}
\]
that is compatible under every adjacent restriction map.
\end{definition}

\begin{definition}[Oracle-valued semantic predicate]
An oracle-valued semantic predicate on a carrier $A$ is a map
\[
  \varphi_s : A \to \TO(\Bool)
\]
where
\[
  \TO(B) = B \times \mathsf{TrustLevel} \times \mathsf{Confidence} \times \mathsf{Provenance}.
\]
\end{definition}

\begin{theorem}[Weakest-link trust composition]
If semantic computations are composed by monadic bind in $\TO$, then the trust of the
composite is bounded above by the meet of the trust levels of the constituents.
\end{theorem}

\begin{proof}[Proof sketch]
The monadic structure is defined to propagate trust via meet. Therefore every downstream
semantic result inherits the weakest trust level seen along the path. This is exactly the
engineering behavior desired by the repository, since high-trust stages should not erase
low-trust inputs.
\end{proof}

\begin{theorem}[Obstruction criterion for hybrid coherence]
If $\Hc{1}(\Layer, \HybridTy) = 0$, then every adjacent pair of layer sections glues to a
coherent hybrid type over the chosen cover.
\end{theorem}

\begin{remark}
This theorem should always be read relative to a chosen cover and chosen restriction maps.
The repository's mature position is strongest when it makes such relativity explicit.
\end{remark}

\begin{proposition}[Mixed-mode compilation as a bridge]
Every mixed-mode surface fragment can be understood as producing either a plain type-level
obligation, a hybrid trust-bearing contract, or both. In particular, user-facing syntax is
not a parallel semantics but a bridge into the existing object language and layered
semantics.
\end{proposition}

\chapter{Appendix: Repository Map}

\begin{longtable}{p{0.28\textwidth}p{0.64\textwidth}}
\toprule
\textbf{Path or namespace} & \textbf{Role in the cohomological typing story} \\
\midrule
\endhead
\repo{src/deppy/types/base.py} & Foundational carrier hierarchy for plain typing. \\
\repo{src/deppy/types/refinement.py} & Refinement predicates and the type-level shape of local facts. \\
\repo{src/deppy/types/dependent.py} & Dependent function and pair formers, identity-like objects, indexed families. \\
\repo{src/deppy/types/subtyping.py} & Ordering and obligation structure for the refinement lattice. \\
\repo{src/deppy/hybrid/core/types.py} & Hybrid type semantics, layer order, trust-facing hybrid type formers. \\
\repo{src/deppy/hybrid/core/checker.py} & Bidirectional hybrid checking and semantic-cache logic. \\
\repo{src/deppy/hybrid/core/algebraic_geometry.py} & Layered presheaves, Cech complexes, cohomological obstruction machinery. \\
\repo{src/deppy/hybrid/core/trust.py} & Trust lattice, evidence classification, provenance plumbing. \\
\repo{src/deppy/hybrid/mixed_mode/syntax.py} & Gracefully degrading user-facing syntax for hybrid programming. \\
\repo{src/deppy/hybrid/discharge/cascade.py} & Staged obligation discharge from trivial proofs to residuals. \\
\repo{src/deppy/hybrid/lean_translation/} & Translation of typed and proof-bearing artifacts into Lean-facing forms. \\
\repo{src/deppy/hybrid/type_expansion/} & Elaboration of prompt or intent material into typed plans and covers. \\
\repo{src/deppy/hybrid/anti_hallucination/} & Artifact typing and structural / semantic hallucination classification. \\
\repo{src/deppy/hybrid/zero_change/} & Inference over existing Python without annotation-first workflows. \\
\repo{src/deppy/equivalence/} & Local-to-global equivalence, semantic diffing, and descent-style reasoning. \\
\repo{src/deppy/render/} & Turning typed results, residuals, and trust labels into readable surfaces. \\
\repo{src/deppy/theory_packs/} & Domain-local solver geometries and library-facing semantics. \\
\repo{agent/orchestrator.py} & Prompt-driven cover growth, generation, verification, and assembly. \\
\repo{web/} & Human-facing explanation layer summarizing the mathematics and APIs. \\
\bottomrule
\end{longtable}

This table should not be read as a dependency graph with one single topological order. It
is a semantic map. The repository's pieces feed one another in multiple directions, but the
map makes clear that the plain and hybrid type systems remain central reference points for
most of the interesting work.


\end{document}
